\documentclass[10pt,a4paper]{amsart}
\usepackage[margin=19mm]{geometry}
\usepackage{amsmath,amssymb,mathtools}
\usepackage[scr=rsfs,cal=euler]{mathalfa}
\usepackage{xfrac}
\usepackage[unicode,hidelinks,bookmarks=false]{hyperref}
\newcommand{\ie}{i.e.\ }
\newcommand{\cf}{cf.\ }
\newcommand{\resp}{respectively\ }
\newcommand{\Subsection}{\subsection}
\providecommand{\nobreakdash}{\nobreak\hbox{-}}
\setlength{\emergencystretch}{2em}
\multlinegap0pt
\usepackage{amssymb}
\usepackage{xcolor}
\usepackage{tikz}
\usepackage[shortlabels]{enumitem}
\newtheorem{proposition}{Proposition}
\newcommand{\biconm}[1]{{\{#1\}''}}
\newcommand{\exte}{\textnormal{ext}}
\newcommand{\inte}{\textnormal{int}}
\newcommand{\ran}{\textnormal{ran}}
\title{Spectral cuts and \\ unconventional functional calculi}
\author{Eva A. Gallardo-Guti\'errez, F. Javier Gonz\'alez-Do\~na}
\date{Source: arXiv:2603.21156v1 (CC BY 4.0)}
\begin{document}
\maketitle

\noindent\emph{Source and licence.} Spectral cuts and \\ unconventional functional calculi. Version arXiv:2603.21156v1 (CC BY 4.0), distributed by arXiv under the Creative Commons Attribution 4.0 International licence, http://creativecommons.org/licenses/by/4.0/, as recorded in the arXiv licence metadata for this exact version and shown on the abstract page https://arxiv.org/abs/2603.21156v1. Full licence notice: \texttt{licenses/arxiv-2603.21156-CC-BY-4.0.txt}.\par\smallskip

\noindent\textit{Editorial scope.} Counted unit: the article's proposition cortan (source0.tex lines 764--779). The article's complete original proof of it is delivered with it (source0.tex lines 782--832, one \texttt{proof} environment of 51 lines). No mathematics is written, repaired or re-derived: the statement and the proof are the article's own lines, in the article's own notation, and no retained line is substituted or re-flowed.  Retained uncounted as the source-local context the proof needs: lines 666--672.  Nothing was omitted from any retained span, so the package carries no omission record.  No retained line contains a citation, so the wrapper carries no bibliography and no bibliography entry.  No retained line contains a figure environment, an image-inclusion command or drawing code, so no image is delivered and the package declares no asset dependency.  Editorial formatting only, in addition: an AMS \texttt{amsart} preamble that restates the article's own declarations and macros used by the retained lines, namely the environments \texttt{proposition} on separate counters, together with the standard packages those lines need.  The retained environments print in this document as Proposition 1 in order of appearance, whereas the article prints the same items with the numbering of its own full document.  The complete native source of the article as distributed in the arXiv e-print for this exact version is bundled unchanged as \texttt{sources/196919/source0.tex}.
\begin{proposition}\label{proposicion union curvas}
Let $T$ be a linear bounded operator on $X$ with the SVEP and $\Gamma_1,\Gamma_2$ two admissible cycles. Assume that $T$ admits plain spectral cuts along both $\Gamma_1$ and $\Gamma_2$ with associated projections $P_{\Gamma_1}$ and $P_{\Gamma_2}$, respectively. Assume further that $\overline{\inte(\Gamma_1)}\cap \overline{\inte(\Gamma_2)}=\varnothing.$ Then:
    \begin{enumerate}
        \item [(i)] The cycle $\Gamma = \Gamma_1\cup \Gamma_2$ can be oriented to form an admissible cycle satisfying $\inte(\Gamma) = \inte(\Gamma_1)\cup \inte(\Gamma_2).$
        \item [(ii)] $T$ admits a plain spectral cut along $\Gamma$, with associated projection $P_\Gamma = P_{\Gamma_1}+P_{\Gamma_2}$.
    \end{enumerate}
\end{proposition}

\begin{proposition}\label{proposicion union curvas que cortan}
Let $T$ be a bounded linear operator on $X$ with the SVEP, and let $\gamma_1,\ldots,\gamma_n$ be positively oriented rectifiable Jordan curves such that $T$ admits a plain spectral cut along each $\gamma_k$, with associated projections $P_{\gamma_k}$. Assume that
$\mathrm{int}(\gamma_i)\cap \mathrm{int}(\gamma_j)=\varnothing \text{ for } i\neq j,
$ and denote $\gamma_{i,j}=\gamma_i\cap\gamma_j$. If
$$
\beta=\overline{(\bigcup\limits_{k=1}^n \gamma_k)\setminus \left(\bigcup_{i,j=1, \ i\neq j}^n \gamma_{i,j} \right) },
$$
then
$$
X = X_T\!\left(\overline{\mathrm{int}(\beta)}\right)
\oplus
X_T\!\left(\overline{\mathrm{ext}(\beta)} \cup
\left(\bigcup_{i\neq j}\gamma_{i,j}\right)\right),
$$
and the associated projection is $P_{\beta}=\sum_{k=1}^n P_{\gamma_k}$.
\end{proposition}

\begin{proof}
Proceeding as in the proof of Proposition \ref{proposicion union curvas}, we first show that $P_\beta \in \biconm{T}$ is an idempotent. It suffices to show that $P_{\gamma_i}P_{\gamma_j}=0$ for all $i\neq j$. Note that
$$\ran(P_{\gamma_i}P_{\gamma_j}) \subseteq   X_T(\overline{\inte(\gamma_i)})\cap X_T(\overline{\inte(\gamma_j)}) = X_T(\gamma_i\cap \gamma_j) \subseteq   X_T(\gamma_i) = \{0\},$$ so $P_\beta$ is an idempotent operator.

\smallskip

Now, let us prove that $\ran(P_\beta) = X_T(\overline{\inte(\beta)})$ and $\ker(P_\beta) = X_T\left(\overline{\exte(\beta)}\cup \left(\cup_{i,j=1, \ i\neq j}^n \gamma_{i,j} \right) \right)$, which will yield the statement. For the range identity, note that
    $$
    X_T(\overline{\inte(\gamma_1)})+\cdots + X_T(\overline{\inte(\gamma_n)})\subseteq   X_T(\overline{\inte(\beta)}),
    $$ so we deduce that $\ran(P_\beta)\subseteq   X_T(\overline{\inte(\beta)})$. Now, let $x\in X_T(\overline{\inte(\beta)})$. Observe that
    $$
    \overline{\inte(\beta)} = \bigcup_{k=1}^n \overline{\inte(\gamma_k)}.
    $$
Write $x=P_{\gamma_1}x+(I-P_{\gamma_1})x$. Since $(I-P_{\gamma_1})x\in
\ker(P_{\gamma_1})=X_T(\overline{\mathrm{ext}(\gamma_1)})$, the hyperinvariance of the spectral subspaces yields
$$
(I-P_{\gamma_1})x\in X_T(\overline{\mathrm{int}(\beta)})\cap
X_T(\overline{\mathrm{ext}(\gamma_1)})
\subseteq X_T\!\left(\bigcup_{k=2}^n \overline{\mathrm{int}(\gamma_k)}\right).
$$
Hence $x=y_1+z_1$, where $y_1\in X_T(\overline{\mathrm{int}(\gamma_1)})$ and
$z_1\in X_T\!\left(\bigcup_{k=2}^n \overline{\mathrm{int}(\gamma_k)}\right)$.

\noindent Proceeding inductively, we obtain
$$
x=\sum_{k=1}^n y_k,
\qquad
y_k\in X_T(\overline{\mathrm{int}(\gamma_k)})=\operatorname{ran}(P_{\gamma_k}).
$$
Since $P_{\gamma_i}P_{\gamma_j}=0$ for $i\neq j$, it follows that
$$
P_\beta x
=\left(\sum_{k=1}^n P_{\gamma_k}\right)\sum_{k=1}^n y_k
=\sum_{k=1}^n y_k
=x.
$$
Therefore $x\in \operatorname{ran}(P_\beta)$, and hence
$$
\operatorname{ran}(P_\beta)=X_T(\overline{\mathrm{int}(\beta)}).
$$



Now, to prove that $\ker(P_\beta) = X_T\left(\overline{\exte(\beta)}\cup \left(\cup_{i,j=1, \ i\neq j}^n \gamma_{i,j} \right) \right)$, we argue as in the proof of Proposition \ref{proposicion union curvas} and deduce that
    \begin{equation*}
        \begin{split}
            \ker(P_\beta) &= \ker\left(\sum_{k=1}^n P_{\gamma_k}\right) = \bigcap\limits_{k=1}^n \ker(P_{\gamma_k}) = \bigcap\limits_{k=1}^n X_T(\overline{\exte(\gamma_k)})  = X_T\left(\bigcap\limits_{k=1}^n \overline{\exte(\gamma_k)} \right) \\& = X_T\left(\overline{\exte(\beta)}\cup \left(\cup_{i,j=1, \ i\neq j}^n \gamma_{i,j} \right) \right).
        \end{split}
    \end{equation*}
which is the desired conclusion.
\end{proof}

\end{document}
